\section{Fairness Definitions}

\begin{frame}{Turning ``Fair'' into Something You Can Measure}

    \begin{block}{Notation for this section}
        \begin{itemize}
            \item $A$ --- the group attribute (e.g.\ $A \in \{\GrpA, \GrpB\}$).
            \item $Y \in \{0,1\}$ --- the true outcome (repaid the loan, re-offended, was ill).
            \item $R \in [0,1]$ --- the score the model outputs.
            \item $\Yhat = \mathbf{1}[R > t]$ --- the decision after thresholding.
        \end{itemize}
    \end{block}

    \vspace{0.8em}

    Almost every group-fairness definition is one of three statements:

    \vspace{0.4em}
    \centering
    \small
    \begin{tabular}{lll}
        \toprule
        \textbf{Family} & \textbf{Statement} & \textbf{Plain reading} \\
        \midrule
        Independence & $\Yhat \perp A$ & the decision rate is the same for every group \\
        Separation   & $\Yhat \perp A \mid Y$ & the error rates are the same for every group \\
        Sufficiency  & $Y \perp A \mid R$ & a score means the same thing for every group \\
        \bottomrule
    \end{tabular}

    \vspace{1em}
    \raggedright
    {\scriptsize Framing after S.~Barocas, M.~Hardt and A.~Narayanan, \emph{Fairness and Machine
    Learning: Limitations and Opportunities}, MIT Press, 2023,
    \href{https://fairmlbook.org/}{fairmlbook.org}, Chapter 3.}
\end{frame}

\begin{frame}{The Object Everything Is Computed From}

    \centering
    \textit{One confusion matrix per group. Every definition below is a constraint tying two cells
    of these tables together across groups.}

    \vspace{0.6em}

    \begin{columns}[c]
        \begin{column}{0.48\textwidth}
            \centering \textbf{Group A}
            \begin{tikzpicture}[scale=0.72]
                \draw[fill=raiGreen!15, thick] (0,0)  rectangle (2.4,2.4);
                \draw[fill=raiRed!15, thick]   (2.4,0) rectangle (4.8,2.4);
                \draw[fill=raiRed!15, thick]   (0,-2.4) rectangle (2.4,0);
                \draw[fill=raiGreen!15, thick] (2.4,-2.4) rectangle (4.8,0);
                \node[font=\bfseries] at (1.2,1.2)  {TP};
                \node[font=\bfseries] at (3.6,1.2)  {FN};
                \node[font=\bfseries] at (1.2,-1.2) {FP};
                \node[font=\bfseries] at (3.6,-1.2) {TN};
                \node[font=\scriptsize\bfseries] at (1.2,2.8) {$\Yhat=1$};
                \node[font=\scriptsize\bfseries] at (3.6,2.8) {$\Yhat=0$};
                \node[font=\scriptsize\bfseries, rotate=90] at (-0.45,1.2)  {$Y=1$};
                \node[font=\scriptsize\bfseries, rotate=90] at (-0.45,-1.2) {$Y=0$};
            \end{tikzpicture}
        \end{column}
        \begin{column}{0.48\textwidth}
            \small
            \begin{itemize}
                \setlength\itemsep{0.45em}
                \item \textbf{Selection rate} $= \Pr(\Yhat = 1)$
                \item \textbf{TPR / recall} $= \Pr(\Yhat=1 \mid Y=1)$
                \item \textbf{FPR} $= \Pr(\Yhat=1 \mid Y=0)$
                \item \textbf{FNR} $= 1 - \text{TPR}$
                \item \textbf{PPV / precision} $= \Pr(Y=1 \mid \Yhat=1)$
                \item \textbf{Prevalence} $p = \Pr(Y=1)$
            \end{itemize}
        \end{column}
    \end{columns}

    \vspace{0.8em}
    \begin{block}{}
        \centering
        \small A fairness \emph{audit} is: build this table per group, then compare a chosen cell
        across groups. The hard part is choosing which cell.
    \end{block}
\end{frame}

\begin{frame}{Definition 1: Demographic Parity (Independence)}

    \begin{block}{Demographic parity / statistical parity}
        \[
            \Pr(\Yhat = 1 \mid A = \GrpA) \;=\; \Pr(\Yhat = 1 \mid A = \GrpB)
        \]
        The model selects the same \emph{proportion} of each group.
    \end{block}

    \vspace{0.6em}

    \begin{columns}[T]
        \begin{column}{0.48\textwidth}
            \textbf{When it is the right question}
            \begin{itemize}\small
                \setlength\itemsep{0.3em}
                \item Allocating an opportunity where you do not trust the label
                      (advertising, outreach, shortlisting).
                \item You believe the historical outcome $Y$ is itself contaminated.
                \item Legal context: the US ``four-fifths rule'' in employment screening flags a
                      selection-rate ratio below $0.8$ as evidence worth investigating.
            \end{itemize}
        \end{column}
        \begin{column}{0.48\textwidth}
            \textbf{Where it goes wrong}
            \begin{itemize}\small
                \setlength\itemsep{0.3em}
                \item It ignores $Y$ entirely. A model can satisfy it by selecting qualified people
                      from one group and arbitrary people from the other.
                \item If the true base rates genuinely differ, enforcing it forces errors.
                \item Measured as a \textbf{ratio} (disparate impact) or a \textbf{difference}
                      --- report which.
            \end{itemize}
        \end{column}
    \end{columns}
\end{frame}

\begin{frame}[allowframebreaks]{Definition 2: Equalised Odds (Separation)}

    \begin{block}{Equalised odds --- Hardt, Price and Srebro (2016)}
        \[
            \Pr(\Yhat = 1 \mid Y = y,\, A = \GrpA) \;=\; \Pr(\Yhat = 1 \mid Y = y,\, A = \GrpB)
            \qquad \text{for } y \in \{0,1\}
        \]
        Equal \textbf{TPR} and equal \textbf{FPR} across groups.
    \end{block}

    \vspace{0.5em}

    \begin{block}{Equal opportunity --- the one-sided version}
        Only $\Pr(\Yhat = 1 \mid Y = 1)$ must match: among people who \emph{would} have succeeded,
        each group has the same chance of being selected.
    \end{block}

    \vspace{0.6em}

    \begin{itemize}\small
        \setlength\itemsep{0.35em}
        \item This is the ProPublica criterion: their complaint was unequal FPR and FNR.
        \item Use it when the burden of a \textbf{mistake} is what matters --- a wrongly denied
              loan, a wrongly flagged defendant, a missed diagnosis.
        \item It conditions on $Y$, so it inherits whatever bias is in the label. If $Y$ is
              ``was arrested,'' equalised odds equalises errors against a biased ground truth.
    \end{itemize}

    \vspace{0.3em}
    {\scriptsize M.~Hardt, E.~Price and N.~Srebro, ``Equality of Opportunity in Supervised Learning,''
    NeurIPS 2016, \href{https://arxiv.org/abs/1610.02413v1}{arXiv:1610.02413v1}.}
\end{frame}

\begin{frame}{Definition 3: Calibration and Predictive Parity (Sufficiency)}

    \begin{block}{Calibration within groups}
        \[
            \Pr(Y = 1 \mid R = r,\, A = \GrpA) \;=\; \Pr(Y = 1 \mid R = r,\, A = \GrpB) \;=\; r
        \]
        A score of $0.7$ means a 70\% chance for everyone, whatever their group.
        \textbf{Predictive parity} is the thresholded version: equal PPV across groups.
    \end{block}

    \vspace{0.8em}

    \begin{itemize}
        \setlength\itemsep{0.45em}
        \item This is the criterion the COMPAS vendor pointed to, and it is a genuinely important
              one: without it, a human reading the score is being misled about what it means.
        \item Well-trained probabilistic models are often \emph{approximately} calibrated
              within groups by default --- which is precisely why the conflict below is unavoidable
              rather than exotic.
        \item Check it with a \textbf{reliability diagram per group}, not with a single number.
    \end{itemize}
\end{frame}

\begin{frame}[allowframebreaks]{You Cannot Have All Three}

    \begin{block}{Chouldechova (2017), Equation (2.6)}
        \[
            \text{FPR} \;=\; \frac{p}{1-p}\cdot\frac{1-\text{PPV}}{\text{PPV}}\cdot\bigl(1-\text{FNR}\bigr)
        \]
        with $p$ the prevalence $\Pr(Y=1)$ in that group.
    \end{block}

    \vspace{0.6em}

    Fix PPV to be equal across two groups. If $p_\GrpA \neq p_\GrpB$, the right-hand side differs,
    so FPR and FNR \textbf{cannot both} be equal. Predictive parity and error-rate balance are
    incompatible whenever base rates differ.

    \vspace{0.8em}

    \begin{itemize}\small
        \setlength\itemsep{0.35em}
        \item Kleinberg, Mullainathan and Raghavan (2016) prove the same tension for calibrated
              \emph{scores}: except in degenerate cases (equal base rates, or a perfect predictor),
              calibration and balance for both classes cannot hold together.
        \item Nothing about this depends on the model, the features, or the amount of data. It is
              arithmetic.
    \end{itemize}

    \vspace{0.3em}
    {\scriptsize A.~Chouldechova, ``Fair prediction with disparate impact,'' \emph{Big Data} 5(2), 2017,
    \href{https://arxiv.org/abs/1703.00056v1}{arXiv:1703.00056v1}. J.~Kleinberg, S.~Mullainathan and
    M.~Raghavan, ``Inherent Trade-Offs in the Fair Determination of Risk Scores,'' ITCS 2017,
    \href{https://arxiv.org/abs/1609.05807v2}{arXiv:1609.05807v2}.}
\end{frame}

\begin{frame}{Seeing the Conflict on a Toy Example}

    Two groups, 100 people each. True positive rate in the population: $p_\GrpA = 0.5$,
    $p_\GrpB = 0.3$. A model that is \textbf{equally accurate within each group}
    (TPR $=0.8$, TNR $=0.8$ for both):

    \vspace{0.6em}
    \centering
    \small
    \begin{tabular}{lcccccc}
        \toprule
        & $Y=1$ & $Y=0$ & TP & FP & \textbf{PPV} & \textbf{FPR} \\
        \midrule
        Group A ($p=0.5$) & 50 & 50 & 40 & 10 & $40/50 = \mathbf{0.80}$ & $10/50 = \mathbf{0.20}$ \\
        Group B ($p=0.3$) & 30 & 70 & 24 & 14 & $24/38 = \mathbf{0.63}$ & $14/70 = \mathbf{0.20}$ \\
        \bottomrule
    \end{tabular}

    \vspace{1em}
    \raggedright

    \begin{itemize}\small
        \setlength\itemsep{0.35em}
        \item Error rates are balanced (equal TPR and FPR) --- but a positive prediction means
              something different in each group (0.80 vs 0.63).
        \item Push PPV to match and the error rates separate. There is no threshold that fixes both.
        \item \textbf{Also note:} the selection rate is 50/100 vs 38/100, so demographic parity
              fails as well.
    \end{itemize}
\end{frame}

\begin{frame}{Group Fairness Is Not the Only Framing}

    \begin{columns}[T]
        \begin{column}{0.48\textwidth}
            \textbf{Individual fairness}
            \vspace{0.3em}

            ``Similar individuals should receive similar predictions'' (Dwork et al., 2012).

            \vspace{0.5em}
            \begin{itemize}\small
                \setlength\itemsep{0.3em}
                \item Attractive, and it catches cases group parity misses.
                \item Requires a task-specific \textbf{similarity metric} --- which is the whole
                      problem, moved somewhere else.
            \end{itemize}
        \end{column}
        \begin{column}{0.48\textwidth}
            \textbf{Causal / counterfactual fairness}
            \vspace{0.3em}

            ``Would the decision change if this person had belonged to another group, holding
            everything else fixed?'' (Kusner et al., 2017).

            \vspace{0.5em}
            \begin{itemize}\small
                \setlength\itemsep{0.3em}
                \item Directly expresses what most people mean by discrimination.
                \item Needs a causal graph you are willing to defend, and the attribute is often
                      not manipulable in the required sense.
            \end{itemize}
        \end{column}
    \end{columns}

    \vspace{1em}
    \begin{block}{}
        \centering
        \small For an introductory course: use group metrics because you can compute them, and
        remember that they are a screening tool, not a certificate.
    \end{block}
\end{frame}

\begin{frame}{So Which Definition Do You Use?}

    This is a decision about \textbf{harms}, not a hyperparameter search. Four questions, in order:

    \vspace{0.8em}

    \begin{enumerate}
        \setlength\itemsep{0.55em}
        \item \textbf{Who is affected, and what is the harm?} Being wrongly \emph{denied}
              something (allocative) or being wrongly \emph{flagged} (punitive) point to different
              error rates.
        \item \textbf{Do you trust the label $Y$?} If the label encodes past discrimination,
              definitions that condition on $Y$ inherit it --- lean towards independence.
        \item \textbf{Is the score read by a human?} If yes, calibration within groups is close to
              mandatory, or you are misinforming the decision-maker.
        \item \textbf{What does the law in your jurisdiction require?} Sometimes the answer is
              fixed for you --- and per-group thresholds may be illegal (see the next section).
    \end{enumerate}

    \vspace{0.9em}
    \begin{block}{}
        \centering
        Write the chosen criterion, and the criteria you consciously gave up, into the model card.
        An unstated choice is still a choice.
    \end{block}
\end{frame}
